% Converted from a Google Doc by gdoc-to-overleaf.

\begin{abstract}
We study when an agent that improves its own improvement procedure converges. Our main result gives a sufficient condition in terms of a contraction constant \(\lambda < 1\), and we show the bound is tight.
\end{abstract}

\section{Introduction}\label{sec:intro}

Recursive self-improvement has been discussed since \citet{good1965}.\footnote{Good's paper is often cited as the origin of the intelligence explosion argument.} We make three contributions:

\begin{enumerate}
\item A formal definition of an \textbf{improvement operator}.
\item A convergence theorem (Theorem~\ref{thm:main}).
\item Experiments, summarised in Table~\ref{tab:conv} and Figure~\ref{fig:steps}.
\end{enumerate}

Related ideas appear in several lines of work:
\begin{itemize}
\item Fixed-point theory
  \begin{itemize}
  \item Banach's theorem
  \item Kleene's recursion theorem
  \end{itemize}
\item Meta-learning % Comment: Cite Schmidhuber (1987) here.
\end{itemize}

\hl{TODO: add a sentence on bounded rationality here.}

\section{Main Result}\label{sec:main}

\begin{definition}
An \emph{improvement operator} is a map \(T:A \rightarrow A\) on a complete metric space of agents.
\end{definition}

\begin{theorem}\label{thm:main}
If \(T\) is a contraction with constant \(\lambda < 1\), then the sequence \(a_{n + 1} = T(a_{n})\) converges to a unique fixed point.
\end{theorem}

\begin{proof}
This is the Banach fixed-point theorem applied to \(T\). The error after \(n\) steps is at most
\[
d(a_{n},a^{*}) \leq \lambda^{n}d(a_{0},a^{*})
\]
which goes to zero.
\end{proof}

\subsection{Implementation}\label{sec:impl}

The update loop is short:
\begin{verbatim}
for step in range(n):
    agent = improve(agent)
\end{verbatim}

\section{Experiments}\label{sec:exp}

Table~\ref{tab:conv} reports convergence speed. We set \(\epsilon = 10^{-3}\) throughout.

\begin{table}[t]
\centering
\caption{Convergence speed on the toy benchmark.}\label{tab:conv}
\begin{tabular}{lrr}
\toprule
\textbf{Method} & \textbf{Steps to converge} & \textbf{Final error} \\
\midrule
Baseline & 120 & 0.031 \\
Ours & 45 & 0.004 \\
\bottomrule
\end{tabular}
\end{table}

\begin{figure}[t]
\centering
\includegraphics[width=3.75in,alt={Bar chart of steps to converge for four settings}]{gdoc/figures/image1.png}
\caption{Steps to converge across four settings.}\label{fig:steps}
\end{figure}

As Figure~\ref{fig:steps} shows, the contraction constant predicts speed well.

\begin{thebibliography}{9}
\bibitem[Good(1965)]{good1965} Good, I. J. (1965). Speculations concerning the first ultraintelligent machine. Advances in Computers, 6, 31--88.

\bibitem[Banach(1922)]{banach1922} Banach, S. (1922). Sur les opérations dans les ensembles abstraits. Fundamenta Mathematicae, 3, 133--181.
\end{thebibliography}
